\section{The graded wheel source and exact decoding}
\label{sec:setup}

\subsection{Wheel recursion}

Set $Q_0=1$ and $Q_1=q_1=2$, and recursively define
\begin{equation}
q_{k+1}=\min\{n>q_k:\gcd(n,Q_k)=1\},
\qquad
Q_{k+1}=Q_kq_{k+1}.
\label{eq:wheel-recursion}
\end{equation}
This rule contains no stored prime table.

\begin{lemma}[Prime enumeration]
\label{lem:prime-enumeration}
For every $k\ge1$, $q_k$ is the $k$th rational prime and
$Q_k=\prod_{j=1}^kq_j$.
\end{lemma}

\begin{proof}
The claim is true for $k=1$.  Assume it through $k$.  Every composite
$n$ strictly between $q_k$ and the next prime has a prime divisor at most
$q_k$, hence shares a factor with $Q_k$.  The next prime shares none.
Therefore the minimum in \eqref{eq:wheel-recursion} is precisely the next
prime, and multiplication gives the formula for $Q_{k+1}$.
\end{proof}

The arithmetic fact needed below is consequently modest but exact: the
sequence $(q_k)$ is injective and unbounded.  The dynamical obstruction would
also apply to any other nonrecurrent chronology; primality enters only
through \cref{lem:prime-enumeration}.

\subsection{Strict grading and target maps}

At level $k\ge0$, let
\[
R_k=\{0\le r<Q_k:\gcd(r,Q_k)=1\}.
\]
For $r\in R_k$, join $r$ to every lift
$r+jQ_k\in R_{k+1}$ with $0\le j<q_{k+1}$ except the unique lift divisible
by $q_{k+1}$.  Let $X_k$ be the set of one-sided infinite paths whose first
edge begins in $R_k$.  Deleting the first edge gives the map $\sigma$ below.
This residue graph fixes the wheel-sieve source, although the obstruction
will use only its grading and exact multiplier clock.

Let
\begin{equation}
X=\bigsqcup_{k\ge0}X_k,
\qquad
\sigma(X_k)\subseteq X_{k+1},
\label{eq:grading}
\end{equation}
be the one-sided wheel tail-path system.  Each $X_k$ is nonempty.  Write
$\ell(x)=k$ for $x\in X_k$ and define the exact integer clock and derived
logarithmic value
\begin{equation}
\kappa(x)=q_{\ell(x)+1},
\qquad
\tau(x)=\log\kappa(x).
\label{eq:clock}
\end{equation}
Then $\ell(\sigma^n x)=\ell(x)+n$ and
$\kappa(\sigma^n x)\ne\kappa(x)$ for every $n\ge1$.

\begin{definition}[Equivariant image and exact decoder]
\label{def:decoder}
Let $S:Y\to Y$ be a self-map.  A source-to-target map
$\pi:X\to Y$ is \emph{equivariant} if
\begin{equation}
S\circ\pi=\pi\circ\sigma.
\label{eq:equivariance}
\end{equation}
An \emph{exact autonomous clock decoder} on the direct image is a
single-valued function $d:\pi(X)\to C$ such that
\begin{equation}
d(\pi(x))=a_{\ell(x)}
\label{eq:decoder}
\end{equation}
for a frozen clock sequence $(a_k)$ in $C$.  The wheel choices are
$a_k=q_{k+1}$ or $a_k=\log q_{k+1}$.  If $\pi$ is onto, it is called a
factor map in the set-theoretic category used by the direct-image theorem.
\end{definition}

The word \emph{autonomous} is substantive.  The decoder may read a state,
edge, finite window, infinite window, or entire target configuration, but it
must define one function of the target point.  A rule that also reads the
absolute time, source level, chosen lift, or traversal number is not a
function on the target phase space.  Classical sliding-block-code
assumptions are therefore sufficient but not necessary for our theorem.

For $m\ge1$, write
\[
\Per_m(S)=\{y\in Y:S^my=y\}.
\]
We distinguish the direct image $\pi(X)$ from an orbit closure
$\cl{\pi(X)}$.  The former needs no topology; the latter requires an
explicit topology and inheritance rule.
